The Introduction section is written to /home/sauravg/Desktop/dataset/paper_detection/section_introduction.tex. Raw LaTeX body below:

\section{Introduction}

Legged robots earn their mobility in exactly the settings that trip them up. A quadruped that can step over a curb or cross rubble is also a quadruped that can hook a foot under a loose cable, snag a calf on a taut wire, or drag a trailing line between its legs. We call this failure \emph{leg entanglement}, and on a walking platform such as the Unitree Go2~\cite{unitree_go2} it is both common and dangerous: the affected leg loses its expected swing, the controller keeps commanding a gait that the physical leg can no longer execute, and the robot lurches, stalls, or falls. Detecting the condition early, and knowing which leg is caught, is a prerequisite for any sensible response.

The obvious sensing route, exteroception, is unusually weak for this problem. Cables and wires are thin, dark, and frequently occluded by the robot's own body or by the very legs they wrap around; the interaction happens below the belly, outside a forward-facing camera's useful field of view, and against cluttered ground. Proprioception is the natural alternative. The joint encoders, joint torques, foot-contact forces, and inertial signals that already drive the walking controller carry a direct mechanical trace of a snag, and proprioceptive signals are routinely used for legged locomotion and state estimation~\cite{hwangbo2019learning,lee2020learning,camurri2017contact}. When a leg is held back, the thigh joint fights a load it cannot move: torque climbs while velocity stays near zero. That is a real signature, and it is the one we build on.

The difficulty is that the signature is subtle and easy to confuse. A high thigh torque at low joint velocity also describes a robot standing still and bearing its own weight, so a naive torque rule fires constantly during ordinary stops. The confusion is sharpest for the rigid ``Lock'' stance the Go2 holds immediately after standing up, where every joint is stiff, torques are large, and motion is nil, yet nothing is wrong. A useful detector therefore cannot simply threshold effort; it must read effort in the temporal context of the gait and separate a genuine snag from a benign, effortful pose.

We frame the task as time-series event detection over the proprioceptive stream and ask three questions at every timestep: \emph{whether} any leg is entangled, \emph{which} of the four legs (front-right, front-left, rear-right, rear-left) it is, and \emph{how severe} the entanglement is. The first is a binary alarm, the second is a per-leg attribution needed to act on the right limb, and the third is a bounded severity index that grades a light brush against a hard snag. Figure~\ref{fig:system} shows how a single causal model answers all three from a short window of on-robot signals, published as a state message for downstream consumers on ROS~2~\cite{macenski2022ros2}. Answering the three jointly, rather than with three unrelated detectors, lets the shared representation carry the temporal context that distinguishes entanglement from standing and from the Lock stance.

This paper makes the following contributions:
\begin{itemize}
  \item A proprioception-only formulation of leg entanglement as a multi-task, per-timestep event-detection problem that jointly answers \emph{whether}, \emph{which}, and \emph{how-severe}, using only signals the walking controller already produces.
  \item A compact causal dilated temporal convolutional network with three heads (detection, per-leg attribution, and a severity auxiliary), together with physics-grounded features that encode the ``high torque, little motion'' snag signature and let the model separate entanglement from standing and from the rigid post-stand-up Lock stance.
  \item A physics-grounded, calibrated per-leg severity index that gates a Mahalanobis deviation from the walking distribution by resistive thigh effort and per-leg confidence, reported as a severity index rather than a measured force.
  \item An evaluation under a leakage-safe, whole-recording split with leave-one-recording-out cross-validation, a rule-based baseline under an identical protocol, feature ablations, and a calibration study, alongside a runtime whose numerical output matches the research pipeline and whose per-window latency sits well under the real-time budget on the deployment CPU.
\end{itemize}